% Source: https://artofproblemsolving.com/community/c3h1988702_largest_natural_number
% Stable identity: aops-topic-1988702-q01
Arrange the numbers in increasing order. Since they are distinct natural numbers, the $k$-th smallest is at least $k-1$, so
\[
2n=x_1+\dots+x_n\ge0+1+\dots+(n-1)=\dfrac{n(n-1)}{2}.
\]
Hence $n-1\le4$, so $n\le5$. For $n=5$ the numbers $0,1,2,3,4$ are distinct and have sum $10=2\cdot5$. The largest $n$ is $\boxed{5}$.

(If the numbers had to be positive integers, the same argument would give $\dfrac{n(n+1)}{2}\le2n$, so $n\le3$, attained by $1,2,3$.)
